\section{Calculus of expansions with transported precision}
\label{sec:series-calculus}

Arithmetic, composition and differentiation must transport the uncertainty
of an approximation together with its finite expression. A useful general
form is
\begin{equation}
 F(y)=b(y)+c(y)\left[J(w(y))+\OO(w(y)^P\M(w(y))^D)\right],
 \qquad w(y)\longrightarrow0^+,
 \label{eq:calculus-carrier}
\end{equation}
where $b,c,w$ are exact expressions and $J$ is a finite power--log jet.
The offset $b$ and prefactor $c$ are treated as exact. Thus a cutoff in
the jet describes relative precision with respect to $c$, and the absolute
remainder contains $|c|$. This form includes ordinary forward and inverse expansions, Lambert expansions and exact coordinate substitutions. A transformed observable generally satisfies a different equation from the original inverse; a residual estimate must refer to the equation that the approximation actually represents.

For addition and multiplication, the precision rules of
Section~\ref{sec:scale} are used after expressing both operands in the same
positive coordinate. An exact prefactor may be absorbed into the jet when
it is a finite power--log expression. Distinct prefactors can be aligned
when their ratio lies in that algebra. Independently ordered exponential factors or incompatible coordinates require a larger common asymptotic scale before a single cutoff can be assigned.

\paragraph{Powers and logarithms.}
For a positive monomial leading term $a w^\alpha$, write
$J=a w^\alpha(1+U)$. A constant power and a logarithm reduce to the unit
series of $(1+U)^r$ and $\log(1+U)$. The input uncertainty first becomes
$\OO(w^{P-\alpha}\M^D)$ in $U$. Multiplication by the output prefactor
then transports it to the correct absolute scale. Real noninteger powers
require positivity on the selected branch. These operations retain any
earlier input precision ceiling; requesting a higher output cutoff does
not create unknown coefficients.

\paragraph{Exponentials and exact carriers.}
Suppose
\[
 H=H_{\le0}+H_{>0}+\OO(w^P\M^D),\qquad P>0,
\]
where $H_{\le0}$ contains the nonpositive-weight blocks. Then
\[
 e^H=e^{H_{\le0}}
 \left[e^{H_{>0}}+\OO(w^P\M^D)\right].
\]
Indeed the omitted argument tends to zero and $e^r=1+\OO(r)$ there.
The operation keeps $e^{H_{\le0}}$ as an exact carrier and expands only
the vanishing part. For example $H=w^{-1}+\log w+w$ yields
$w e^{1/w}(1+w+w^2/2+\cdots)$. This also supplies the precision rule
needed when a source logarithm is inverted and the source variable must
be reconstructed by exponentiation. A nonvanishing argument uncertainty cannot in general justify a relative error tending to zero after exponentiation.

\paragraph{Composition and regular observables.}
The outer local coordinate must become $a w^\alpha(1+o(1))$ with
$a,\alpha>0$ after substitution of the inner expansion. This condition
checks both the limiting point and its selected side. An outer remainder
$\OO(v^P\M(v)^D)$ becomes $\OO(w^{\alpha P}\M(w)^D)$, since
$\log v=\alpha\log w+\OO(1)$. The uncertainty of the inner argument is
propagated independently through every jet operation, and the worse
precision is retained. Regular analytic observables at a finite limiting
argument use their Taylor series and the same precision arithmetic.
Poles, branch points and logarithmic leading coefficient functions require
their respective scale normalizations; the regular Taylor-composition rule does not apply at a singular point.

The outer estimate must also hold uniformly along the substituted parameter
path. More precisely, a bound
$|F(v;\lambda)-J(v;\lambda)|\le C(\lambda)v^P\M(v)^D$
for $0<v<\delta(\lambda)$ can be used at
$(v,\lambda)=(v(w),\lambda(w))$ only when this path remains in the stated
domain. Its transported envelope still contains $C(\lambda(w))$;
discarding that factor requires a uniform upper bound. Estimates proved
separately for each fixed parameter do not provide such a bound, as in
Section~\ref{sec:formal-successive-jets}. For example, at each fixed $a>0$,
\[
 \frac{x}{a+x}=\frac{x}{a}
   -\frac{x^2}{a(a+x)}
   =\frac{x}{a}+\OO_a(x^2),\qquad x\longrightarrow0^+.
\]
On the diagonal $x=a\longrightarrow0^+$ the exact value is $1/2$, whereas
the displayed finite term is $1$ and the discarded term is $-1/2$.
The parameter dependence can therefore matter even when it occurs only
in a discarded term. Re-expanding the complete function along the joint
path can supply a new valid estimate. An exact outer identity has no
outer remainder to transport, but its parameter and branch conditions
must still hold, and uncertainty in the inner argument must still be
propagated. Refinement in the original fixed-parameter regime alone
does not establish a uniform estimate on a new path.

\paragraph{Taylor order and sided limiting constants.}
Let $\sigma\in\{-1,1\}$ and suppose the selected germ has the proved estimate
\begin{equation}
 h(c+\sigma v)=\sum_{k=0}^{q-1}a_kv^k+\OO(v^q),
 \qquad v\longrightarrow0^+,
 \label{eq:sided-taylor-contract}
\end{equation}
where $q$ is a positive integer. For example, an analytic extension of
$v\mapsto h(c+\sigma v)$ to a neighborhood of zero supplies this estimate
with a bounded Taylor remainder. A formal coefficient list alone does not.
An absent coefficient below $q$ may denote a known zero; coefficients at
or above $q$ remain undetermined by this estimate. Thus the truthful
description $v-v^3/6=v+\OO(v^2)$ does not imply
$v-v^3/6=v+\OO(v^5)$.

Suppose the complete inner germ satisfies
$V(w)=\sigma(g(w)-c)>0$ eventually and
$V(w)=\OO(w^\alpha\M(w)^d)$, with $\alpha>0$ and $d\ge0$.
For an exclusive power cutoff $C>0$, set $N=\lceil C/\alpha\rceil$.
If $q\ge N$, then
\begin{equation}
 h(g(w))=\sum_{k=0}^{N-1}a_kV(w)^k
   +\OO\!\left(w^{\alpha N}\M(w)^{dN}\right).
 \label{eq:sided-taylor-demand}
\end{equation}
Indeed, the discarded Taylor powers and the remainder in
\eqref{eq:sided-taylor-contract} together are $\OO(V^N)$, since $V\to0$;
the displayed bound follows by taking the $N$th power of its majorant.
All orders $k<N$ have $k\alpha<C$. If $\alpha N=C$, the logarithmic degree
$dN$ at that boundary must be retained. If $\alpha N>C$, the strict power
margin absorbs this fixed logarithmic factor into $\OO(w^C)$.
When $q<N$, the available Taylor remainder instead contributes
$\OO(w^{\alpha q}\M^{dq})$ and cannot be assigned the requested higher
power without further information. In either case, uncertainty in an
approximation to $V$ and blocks discarded from its finite powers contribute
their own bounds; \eqref{eq:sided-taylor-demand} does not remove them.

The constant $a_0$ is the sided limit in
\eqref{eq:sided-taylor-contract}, and need not equal the point value $h(c)$.
For the fractional-part function $h(t)=t-\lfloor t\rfloor$ near $c=1$,
\[
 h(1)=0,\qquad h(1-v)=1-v,\qquad h(1+v)=v
 \quad(0<v<1).
\]
An exactly constant inner germ uses $h(c)$. A proved left or right approach
uses the corresponding sided germ, including its constant; changing only
$h(c)$ cannot affect a composition eventually separated from $c$.
An estimate $g=1+\OO(w)$ alone supplies neither a side nor that separation:
it permits $1-w$, $1+w$, and the constant $1$. A Taylor contract valid on
every permitted approach, including the point if it can be attained, is
required before one such approximation can describe all these possibilities.
The positivity of $\sigma(g-c)$ must therefore follow from the full inner
germ and its error bound, not merely from a chosen formal Taylor coordinate.

\paragraph{Differentiation requires derivative information.}
An estimate $R(w)=\OO(w^P\M(w)^D)$ alone does not imply any estimate on
$R'$. For example $w^P\sin(e^{1/w})$ has that magnitude bound and a much
larger derivative. Therefore a nonzero remainder can be differentiated
$n$ times only with the additional hypotheses
\[
 R^{(j)}(w)=\OO(w^{P-j}\M(w)^D),\qquad 0\le j\le n.
\]
Each differentiation uses one order of these derivative bounds. The jet derivative uses
$D_w[w^\beta Q(\log w)]=w^{\beta-1}(\beta Q+Q')$; the derivatives
of the exact coordinate, offset and prefactor in
\eqref{eq:calculus-carrier} are included by the chain and product rules.
Exact finite expressions require no remainder hypothesis. A discarded tail
of a known finite power--log expression also has derivative bounds of every
order. In contrast, a derivative bound established for one unknown remainder
is not automatically promoted to a stronger exponent when the source is
refined; that stronger bound must be established separately.

\paragraph{Truncation and refinement.}
Truncation moves the first discarded block, including its logarithmic
degree, into the remainder. It cannot improve uncertainty inherited from
the input. Refinement instead determines more coefficients from the original function,
implicit equation or composition. It remains subject to any error bound
inherited from unspecified input data. Changing only the exponent written
in a remainder symbol supplies no new mathematical information.

A remainder of a known analytic function and an error in an unspecified
function play different roles. The Taylor expansion of $\sin x$ can be
continued to arbitrary finite order because the original function is known.
If only a finite approximation $T$ and a bound $F-T=\OO(w^P\M^D)$ are given,
then the unknown coefficients of $F$ at or beyond that precision are not
determined. Higher-order operations on $T$ cannot remove this uncertainty.

\subsection{Separate error envelopes}
\label{sec:separate-error-envelopes}

A common ordered coefficient algebra is unnecessary for transporting
magnitude bounds. Along one fixed real approach, suppose
$F=A+\OO(R)$ and $G=B+\OO(S)$, where $F,G,A,B$ are real-valued and
$R,S\ge0$ are error envelopes. Their coordinates and rates of decay may
differ. The arithmetic of Section~\ref{sec:evaluator} then has the
following form, without selecting a single exponent cutoff.

\begin{proposition}[Error envelopes for compound functions]
\label{prop:separate-error-envelopes}
Under these hypotheses,
\begin{equation}
\begin{aligned}
 F+G&=A+B+\OO(R+S),\\
 FG&=AB+\OO\bigl(|A|S+|B|R+RS\bigr).
\end{aligned}
\label{eq:separate-envelope-arithmetic}
\end{equation}
For every fixed positive integer $n$,
\begin{equation}
 F^n=A^n+\OO\!\left(
   \sum_{k=1}^{n}\binom nk |A|^{n-k}R^k\right),
\label{eq:separate-envelope-integer-power}
\end{equation}
where the factor $|A|^0$ in the last summand is $1$, including at $A=0$.
No nonvanishing or relative-smallness hypothesis is needed here.

If $A$ is eventually nonzero and $R/|A|\to0$, then for fixed real $r$
\[
 F^r=A^r+\OO\bigl(|A|^{r-1}R\bigr),
\]
provided $A>0$ eventually when $r$ is not an integer. In particular,
$F^{-1}=A^{-1}+\OO(R/|A|^2)$. The logarithm and exponential obey
\begin{equation}
\begin{aligned}
 \log F&=\log A+\OO(R/A)
   &&\text{if }A>0\text{ eventually and }R/A\to0,\\
 e^F&=e^A+\OO(e^A R)
   &&\text{if }R\to0.
\end{aligned}
\label{eq:separate-envelope-log-exp}
\end{equation}
Finally $\phi(F)=\phi(A)+\OO(R)$ for
$\phi(t)=|t|,\sin t,\cos t$, without a smallness hypothesis on $R$ or
a finite-limit hypothesis on $A$.
\end{proposition}

\begin{proof}
Write $F=A+u$ and $G=B+v$, with $|u|\le CR$ and $|v|\le DS$
eventually. Addition and the identity
$FG-AB=Av+Bu+uv$ give \eqref{eq:separate-envelope-arithmetic}.
The binomial formula gives \eqref{eq:separate-envelope-integer-power},
absorbing the finitely many factors $C^k$ into one constant.
Under relative smallness, $u/A\to0$, so $F/A$ stays near $1$ and
$F$ has the same eventual sign as $A$. The functions $(1+t)^r$ and
$\log(1+t)$ have bounded derivatives near $t=0$. Apply their mean value
bounds to $t=u/A$ and multiply by the exact prefactor.
If $R\to0$, then $u\to0$ and $e^u-1=\OO(u)$, giving the exponential
bound. Absolute value, sine and cosine are globally $1$-Lipschitz on
the real line, giving the final assertion. As in
Proposition~\ref{prop:soundness}, induction permits finite compositions
whenever each operation satisfies its stated hypotheses.
\end{proof}

These are asymptotic bounds with unspecified constants and thresholds;
the constants may depend on fixed data and powers. Cancellation of
$A+B$ supplies no cancellation of unknown errors without additional
relations between them. The retained expression need not be ordered as
a Poincar\'e expansion, and a trigonometric bound may identify no leading
term at all. Thus these assertions alone provide neither a single
cutoff nor a pointwise numerical certificate, and they imply no
derivative bounds. For a pure error, the integer-power rule gives
$\OO(R)^n=\OO(R^n)$; a fractional power additionally needs a real-branch
condition that a magnitude bound alone does not supply.

\subsection{Magnitude bounds do not transport differentiability}
\label{sec:magnitude-vs-derivative}

The Lipschitz argument above transports a magnitude bound through an absolute
value, and Remark~\ref{rem:modulus-transport} does the same for the modulus of
a complex-valued expansion. Neither transports a \emph{classical derivative}
contract, even when the input is smooth and its recorded derivative bounds are
true. The obstruction is that $t\mapsto|t|$ fails to be differentiable at
$t=0$, and a quantity that changes sign infinitely often near the endpoint
meets that point infinitely often.

\begin{example}[A truthful derivative contract destroyed by an absolute value]
\label{ex:absolute-value-cusps}
Let $f(x)=x+x^2\sin\log x$ for $x>0$, and write $L=\log x$. Then $f$ is
smooth on $(0,\infty)$, $f(x)-x=\OO(x^2)$, and
\[
 \varphi(x):=f'(x)-1=x\,(2\sin L+\cos L)=\sqrt5\,x\,\sin(L+\theta),
 \qquad \tan\theta=\tfrac12 .
\]
Hence $|\varphi|\le\sqrt5\,x$, so $f'(x)>0$ for $x<1/\sqrt5$ and $f$ is
strictly increasing there, with a local inverse $g$ near $0$. Also
$f''(x)=(2\sin L+\cos L)+(2\cos L-\sin L)=\OO(1)$, so the three assertions
\[
 g(y)-y=\OO(y^2),\qquad g'(y)-1=\OO(y),\qquad g''(y)=\OO(1)
\]
are all true, the second because $g'-1=-\varphi(g)/f'(g)$ and $g\sim y$.

Nevertheless $|\varphi|$ is differentiable on no interval $(0,\epsilon)$. Its
zeros in $(0,1)$ are $x_n=\exp(-\arctan\tfrac12-n\pi)$ for $n\ge1$, which
accumulate at $0$, and there $2\sin L+\cos L=0$, so
\[
 \varphi'(x_n)=f''(x_n)=2\cos L-\sin L=\pm\sqrt5\ne0 .
\]
Thus $\varphi$ changes sign at each $x_n$ and $|\varphi|$ has one-sided
derivatives $-\sqrt5$ and $+\sqrt5$ there. The magnitude bound
$|\varphi|=\OO(x)$ survives the absolute value; the derivative bound
$\varphi'=\OO(1)$ does not extend to $|\varphi|$, whose derivative fails to
exist at every $x_n$. The same happens to $|g'(y)-1|$ at the images
$y_n=f(x_n)$, because $g'(y)-1$ vanishes exactly where $f'(g(y))=1$ and
$f''\ne0$ at those points.
\end{example}

A derivative contract must therefore be recorded separately from a magnitude
contract, and lowered whenever an absolute value or a modulus is applied to a
quantity whose sign is not proved eventually constant. A nonvanishing proof on
the relevant neighborhood, or an exact factorization that exhibits the sign,
restores the higher contract; a Lipschitz constant alone does not. Neither does
the accumulation of cusps contradict the magnitude estimate: the two contracts
are independent, and only the weaker one is preserved.
